\begin{resumen}
La generación automática de consultas a bases de datos a partir de lenguaje natural (\emph{text-to-query}) ha avanzado de la mano de los modelos extensos de lenguaje, pero en entornos empresariales aún produce fallas sistemáticas de fidelidad estructural: referencias a tablas o columnas inexistentes, uniones inválidas y patrones de grafo incoherentes con el esquema. El problema se intensifica bajo el estándar híbrido \ac{SQL}/\ac{PGQ} (ISO/IEC 9075-16:2023), donde la generación debe operar simultáneamente sobre un esquema relacional y un grafo de propiedad declarado sobre este.

Esta tesis propone una arquitectura de razonamiento estructural basada en una representación intermedia tipada, IR-SQL/PGQ, que media entre la salida del modelo de lenguaje y la consulta ejecutable. La arquitectura articula cuatro componentes: un anclaje al esquema dual que reduce el espacio de búsqueda; la representación intermedia, expresable de forma uniforme sobre los regímenes relacional, de grafo e híbrido y decidible mediante reglas locales de tipado; un compilador determinista hacia \ac{SQL}/\ac{PGQ}; y un bucle de verificación en cascada que opera sobre la propia representación intermedia con retroalimentación estructurada por descriptores categóricos.

La validación experimental se desarrollará sobre los \emph{benchmarks} Spider y BIRD para el régimen relacional y sobre un corpus propio para el régimen híbrido \ac{SQL}/\ac{PGQ}. Las métricas reportadas incluirán corrección, tasa de alucinación, transferibilidad a entornos heterogéneos y coste agregado por consulta. La hipótesis central, sujeta a verificación empírica, es que la mediación por representación intermedia tipada con verificación en cascada reduce las alucinaciones referenciales y estructurales respecto a la generación directa, sin comprometer el coste sostenido en escenarios empresariales.
\end{resumen}
